% ============================================================
%  ch03a：§3 谱 Eisenstein 级数：引言、3.1 自守 Eisenstein 级数、
%         3.2 一些一般构造（原文 pp.16–19）
% ============================================================
\section{谱 Eisenstein 级数}

我们现在开始走向定理 A。本文余下部分，我们始终假定 $k = \Fb$。

本节的目标是定义并研究如下映射：
\[
  \Eisspec : \Gamma(\LS^{\arthm}_{\bT}, \omega_{\LS^{\arthm}_{\bT}}) \longrightarrow
  \Gamma(\LS^{\arthm}_{\bG}, \omega_{\LS^{\arthm}_{\bG}}).
\]

本节通篇只考虑局部几乎有限型的（预）叠；我们不再一一提及这一假设。

\subsection{自守 Eisenstein 级数}

我们先回顾关于几何 Eisenstein 级数及其函数论对应物——伪 Eisenstein 级数——的若干构造。稍后我们希望在谱侧找到这些构造的对应物。

\pnum{首先，我们有一个典范函子：}
\[
  \Eis_! : \Shv(\Bun_T) \longrightarrow \Shv(\Bun_G),
\]
它由沿 $\Bun_B \to \Bun_T$ 的 $*$-拉回、再沿 $\Bun_B \to \Bun_G$ 的 $!$-推前定义。

由 [6] 中建立的 $\Eis_!$ 的 Hecke 性质，以及 [2] 定理 14.4.3（“Nadler--Yun 定理的逆”，参看 §2.5.3），我们得到：

\begin{prop}
函子 $\Eis_!$ 把 $\qLisse(\Bun_T)$（$= \Shv_{\Nilp_T}(\Bun_T)$）映入 $\Shv_{\Nilp}(\Bun_G)$。
\end{prop}

\pnumt{伪 Eisenstein 级数。}{}\ 令 $\text{ps-Eis} : \Aut^{\unr}_{T,c} \to \Aut^{\unr}_{G,c}$ 为伪 Eisenstein 映射。按定义，它是如下复合：
\[
  \Aut^{\unr}_{T,c} := \mathrm{Func}(\Bun_T(\F_q)) \longrightarrow
  \mathrm{Func}(\Bun_B(\F_q)) \longrightarrow \mathrm{Func}(\Bun_G(\F_q)) := \Aut^{\unr}_{G,c},
\]
即先作限制（注意 $\Bun_B(\F_q) \to \Bun_T(\F_q)$ 的纤维有限），再沿 $\Bun_B(\F_q) \to \Bun_G(\F_q)$ 的纤维求和（这是良定义的，因为我们考虑的是具有有限支撑的函数）。

\pnumt{两者的相容性。}{}\ 函子 $\Eis_!$ 显然缠绕 Frobenius 并保持紧对象，因此我们可以过渡到 Frobenius 的迹，得到映射：
\[
  \tr(\Eis_!) : \tr_{\qLisse(\Bun_T)}(\Frob) \longrightarrow \tr_{\Shv_{\Nilp}(\Bun_G)}(\Frob).
\]
由 [4] 的主定理，我们有同构：
\begin{equation}
  \begin{aligned}
    \tr_{\qLisse(\Bun_T)}(\Frob) &\cong \Aut^{\unr}_{T,c},\\
    \tr_{\Shv_{\Nilp}(\Bun_G)}(\Frob) &\cong \Aut^{\unr}_{G,c}.
  \end{aligned}
\end{equation}

于是 $\tr(\Eis_!)$ 对应于映射 $\Aut^{\unr}_{T,c} \to \Aut^{\unr}_{G,c}$。

由 [4] 定理 5.2.3\footnote{我们指出，前引文献（loc.\ cit.）的陈述（甚至其表述本身）都以 $\Shv_{\Nilp}(\Bun_G)$ 上某个技术性假设为条件；见 [4] §5.1。这一假设最近由作者与 Gaitsgory 验证，将出现在即将发表的工作中。}，同构 (3.1) 由层—函数对应的一个版本给出；由此形式上可得交换图：
\begin{equation}
  \begin{tikzcd}[column sep=large, row sep=large]
    \tr_{\qLisse(\Bun_T)}(\Frob) \arrow[r, "\tr(\Eis_!)"] \arrow[d, "\cong"']
      & \tr_{\Shv_{\Nilp}(\Bun_G)}(\Frob) \arrow[d, "\cong"'] \\
    \Aut^{\unr}_{T,c} \arrow[r, "\text{ps-Eis}"]
      & \Aut^{\unr}_{G,c}.
  \end{tikzcd}
\end{equation}

换言之：\emph{几何 Eisenstein 级数的迹就是伪 Eisenstein 级数。}

\subsection{一些一般构造}

以下除非另有说明，我们只在域 $\e$ 上工作。

下面，我们在一般的叠论语境中给出 $\Eisspec$ 的一个一般构造。

\pnumt{叠的术语。}{}\ 回忆 [8] 中 QCA 叠\zy{QCA 是“quasi-compact with affine diagonal”（拟紧且对角仿射）的缩写；这一概念由 [8] 引入并在其中系统使用。[来源：A/原文 §3.2.1 及其所引 [8]]}这一技术性概念：该术语指拟紧且对角仿射的代数叠。对任意这样的 QCA 叠 $\mathcal{Y}$，[8] 定理 0.4.5 断言 $\IndCoh(\mathcal{Y}) = \mathrm{Ind}(\Coh(\mathcal{Y}))$。此外，由 [8] §3，QCA 叠上的归纳凝聚层有良好的推前理论。

我们也回忆 [2] §5.2 中 ind-代数叠的概念。我们提醒：预叠 $\mathcal{Y}$ 称为\emph{ind-代数的}，如果它是收敛的，且对每个 $n \ge 0$，它的 $n$-截断 ${}^{\le n}\mathcal{Y}$ 可以写成若干 $n$-截断代数叠 $\mathcal{Y}_i$ 沿闭嵌入的滤余极限。如果这些 $\mathcal{Y}_i$ 还可以取为 QCA 的，我们就说 $\mathcal{Y}$ 是\emph{ind-QCA} 的。

我们的主要例子是仿射代数群 $H$ 的 $\LS^{\restr}_H$。由 [2] 推论 5.2.6，$\LS^{\restr}_H$ 是 ind-代数的；此外，该结果的证明表明 $\LS^{\restr}_H$ 实际上是 ind-QCA 的。

由上述 Drinfeld--Gaitsgory 定理，任意 ind-QCA 叠 $\mathcal{Y}$ 的 $\IndCoh(\mathcal{Y})$ 都是紧生成的。同样，ind-QCA 叠之间的态射有良好的 $\IndCoh$-推前理论。

\pnum{下面，我们固定 ind-QCA 代数叠之间的一个 1-可表\footnote{即纤维是代数叠。确切地说，对每个 $S \in \mathrm{AffSch}$ 与 $S \to \mathcal{Z}$，纤维积 $\mathcal{Y} \times_{\mathcal{Z}} S$ 是代数叠。在我们的语境中，这一条件排除了像 $\mathcal{Y} \to \Spec(\e)$ 这样的映射，除非 $\mathcal{Y}$ 是真正的（而非 ind-）代数叠。}映射 $f : \mathcal{Y} \to \mathcal{Z}$。}

此外假设给定自同构 $\phi_{\mathcal{Y}} : \mathcal{Y} \to \mathcal{Y}$ 与 $\phi_{\mathcal{Z}} : \mathcal{Z} \to \mathcal{Z}$，它们被 $f$ 缠绕（即给定同构 $\phi_{\mathcal{Z}} \circ f \simeq f \circ \phi_{\mathcal{Y}}$）。有时我们省略下标，对 $\phi_{\mathcal{Y}}$ 与 $\phi_{\mathcal{Z}}$ 都简写为 $\phi$。

我们构造 $\phi$ 的不动点叠 $\mathcal{Y}^\phi$（相应地 $\mathcal{Z}^\phi$）。确切地说，它是等化子 $\mathrm{Eq}(\mathcal{Y} \rightrightarrows \mathcal{Y})$（两个映射为 $\phi$ 与 $\mathrm{id}$），也可以写成 $\mathcal{Y} \xrightarrow{\mathrm{Graph}_\phi} \mathcal{Y} \times \mathcal{Y} \xleftarrow{\Delta} \mathcal{Y}$ 的纤维积。由假设，我们有诱导映射 $\mathcal{Y}^\phi \to \mathcal{Z}^\phi$，记作 $f^\phi$。

下面，在适当的假设下，我们将构造 $\Gamma(\mathcal{Y}^\phi, \omega_{\mathcal{Y}^\phi})$ 与 $\Gamma(\mathcal{Z}^\phi, \omega_{\mathcal{Z}^\phi})$ 之间的典范映射。

我们把讨论组织如下。在 §3.2.3 与 §3.2.4 中，我们利用 $\IndCoh$ 的标准函子性给出这些映射的“初等”构造。后一个构造尤其有些复杂。之后，我们解释一个使这些构造的存在性变得显然的概念框架（迹的函子性）。作者本人是从后一种观点来思考这些映射的，但担心若缺少具体描述，读者不会相信这套魔法；不需要这种说服的读者可以跳过 §3.2.3--3.2.4。

我们还要指出，ind-QCA 假设是过强的。§3.2.3 并不需要它。它在 §3.2.4 中被轻度使用，用于各种推前函子的存在性，不过更弱的假设已经足够。从根本上说，从更概念化的迹的视角看，这一假设是自然的，它保证 $\IndCoh$ 的紧生成性（从而保证可对偶性）。

\pnumt{推前。}{}\ 首先，假设映射 $f$ 可表且紧合（proper）。此时我们将构造映射：
\begin{equation}
  \Gamma^{\IndCoh}(\mathcal{Y}^\phi, \omega_{\mathcal{Y}^\phi}) \longrightarrow
  \Gamma^{\IndCoh}(\mathcal{Z}^\phi, \omega_{\mathcal{Z}^\phi}).
\end{equation}

事实上，这相当容易。此时映射 $f^\phi : \mathcal{Y}^\phi \to \mathcal{Z}^\phi$ 也是紧合（且可表）的\footnote{确实，由于 $f$ 可表且分离，态射 $\mathcal{Y}^\phi = \mathcal{Y} \times_{\mathcal{Y} \times \mathcal{Y}} \mathcal{Y} \to \mathcal{Y} \times_{\mathcal{Z} \times \mathcal{Z}} \mathcal{Z}$ 是闭嵌入。显然，进一步的投影 $\mathcal{Y} \times_{\mathcal{Z} \times \mathcal{Z}} \mathcal{Z} \to \mathcal{Z} \times_{\mathcal{Z} \times \mathcal{Z}} \mathcal{Z} = \mathcal{Z}^\phi$ 是紧合的。}，而这正是下面所需的全部。

于是 $f^{\phi,\IndCoh}_*$ 左伴随于 $f^{\phi,!}$，故我们得到典范伴随映射 $f^{\phi,\IndCoh}_* f^{\phi,!} \to \mathrm{id}$。这给出映射 $f^{\phi,\IndCoh}_*(\omega_{\mathcal{Y}^\phi}) \to \omega_{\mathcal{Z}^\phi}$；再作用 $\Gamma^{\IndCoh}(\mathcal{Z}^\phi, -)$，即得所需映射 (3.3)。

\begin{remark}
映射 (3.3) 是 $\Gamma(\mathcal{Z}^\phi, \OO_{\mathcal{Z}^\phi})$-线性的。
\end{remark}

\pnumt{拉回。}{}\ 其次，假设 $f$ 最终余连通（且 1-可表）。此时我们将构造映射：
\begin{equation}
  \Gamma^{\IndCoh}(\mathcal{Z}^\phi, \omega_{\mathcal{Z}^\phi}) \longrightarrow
  \Gamma^{\IndCoh}(\mathcal{Y}^\phi, \omega_{\mathcal{Y}^\phi}).
\end{equation}

首先注意\footnote{这一点以及本节中其他类似断言，形式上都可以归约到 [17] 第 4 章 §3 的结果。}，此时 $f^{!,\IndCoh} : \IndCoh(\mathcal{Y}) \to \IndCoh(\mathcal{Z})$ 具有左伴随 $f^{*,\IndCoh}$。我们有 $*$-推前与 $*$-拉回之间的基变换（后者只对最终余连通的态射考虑）。

此时，我们有自然变换
\[
  \begin{aligned}
    &(\mathrm{id} \times f)^{*,\IndCoh}(\mathrm{id} \times \phi_{\mathcal{Z}})^{\IndCoh}_*
      (\mathrm{id} \times f)^{\IndCoh}_* = {}\\
    &(\mathrm{id} \times f)^{*,\IndCoh}(\mathrm{id} \times f)^{\IndCoh}_*
      (\mathrm{id} \times \phi_{\mathcal{Y}})^{\IndCoh}_*\\
    &\qquad\qquad \longrightarrow (\mathrm{id} \times \phi_{\mathcal{Y}})^{\IndCoh}_*
      \in \mathrm{End}(\IndCoh(\mathcal{Y} \times \mathcal{Y})),
  \end{aligned}
\]
它来自伴随。把这一映射作用于 $\Delta^{\IndCoh}_*(\omega_{\mathcal{Y}})$，我们得到典范映射：
\begin{equation}
  \begin{aligned}
    &(\mathrm{id} \times f)^{*,\IndCoh}
      \mathrm{Graph}^{\IndCoh}_{\phi_{\mathcal{Z}} \circ f,*}(\omega_{\mathcal{Y}})\\
    &\qquad\qquad \longrightarrow \mathrm{Graph}_{\phi_{\mathcal{Y}},*}(\omega_{\mathcal{Y}})
      \in \IndCoh(\mathcal{Y} \times \mathcal{Y}).
  \end{aligned}
\end{equation}

这里，对映射 $g : S \to T$，映射 $\mathrm{Graph}_g : S \to S \times T$ 是图像态射，即 $\mathrm{Graph}_g := (\mathrm{id} \times g) \circ \Delta_S$。

令 $\varpi_{\mathcal{Y}}$ 表示典范映射 $\mathcal{Y}^\phi \to \mathcal{Y}$，它把一对 $(y \in \mathcal{Y}, y \simeq \phi_{\mathcal{Z}}(y))$ 送到 $y$；$\varpi_{\mathcal{Z}}$ 类似。% [Tutor Note: 原文此处作 φ_𝒵(y)，但 𝒴^φ 由 φ_𝒴 定义，语义上应为 φ_𝒴(y)；译文保留原文写法，见报告。]

下面，我们将构造典范同构：
\begin{equation}
  \begin{aligned}
    &\Delta^!_{\mathcal{Y}}(\mathrm{id} \times f)^{*,\IndCoh}
      \mathrm{Graph}^{\IndCoh}_{\phi_{\mathcal{Z}} \circ f,*}(\omega_{\mathcal{Y}})\\
    &\qquad\qquad \simeq f^{*,\IndCoh}\varpi^{\IndCoh}_{\mathcal{Z},*}(\omega_{\mathcal{Z}^\phi})
      \in \IndCoh(\mathcal{Y}).
  \end{aligned}
\end{equation}

假设这一构造已经给出，则把 $\Delta^!_{\mathcal{Y}}$ 作用于 (3.5)，我们得到典范映射
\[
  f^{*,\IndCoh}\varpi^{\IndCoh}_{\mathcal{Z},*}(\omega_{\mathcal{Z}^\phi})
  \longrightarrow \Delta^!_{\mathcal{Y}}\mathrm{Graph}_{\phi_{\mathcal{Y}},*}(\omega_{\mathcal{Y}})
  \in \IndCoh(\mathcal{Y}).
\]
由伴随，这又给出典范映射
\[
  \begin{aligned}
    \varpi^{\IndCoh}_{\mathcal{Z},*}(\omega_{\mathcal{Z}^\phi})
    &\longrightarrow f^{\IndCoh}_*\Delta^!_{\mathcal{Y}}
      \mathrm{Graph}_{\phi_{\mathcal{Y}},*}(\omega_{\mathcal{Y}})\\
    &\simeq f^{\IndCoh}_*\varpi^{\IndCoh}_{\mathcal{Y},*}(\omega_{\mathcal{Y}^\phi})
      \in \IndCoh(\mathcal{Z}).
  \end{aligned}
\]
这里我们用了基变换同构 $\Delta^!_{\mathcal{Y}}\mathrm{Graph}^{\IndCoh}_{\phi_{\mathcal{Y}},*} \simeq \varpi^{\IndCoh}_{\mathcal{Y},*}\varpi^!_{\mathcal{Y}}$。现在把整体截面函子 $\Gamma^{\IndCoh}(\mathcal{Z}, -)$ 作用于上式两端，即得所需映射 (3.4)。

还剩下同构 (3.6) 的给出。首先，纤维积图表
\[
  \begin{tikzcd}[column sep=large, row sep=large]
    \mathcal{Y} \arrow[r, "\mathrm{Graph}_{\phi_{\mathcal{Z}} \circ f}"] \arrow[d, "f"]
      & \mathcal{Y} \times \mathcal{Z} \arrow[d, "f \times \mathrm{id}"] \\
    \mathcal{Z} \arrow[r, "\mathrm{Graph}_{\phi_{\mathcal{Z}}}"]
      & \mathcal{Z} \times \mathcal{Z}
  \end{tikzcd}
\]
给出基变换同构
\[
  \mathrm{Graph}^{\IndCoh}_{\phi_{\mathcal{Z}} \circ f,*}(\omega_{\mathcal{Y}})
  \simeq (f \times \mathrm{id})^!\mathrm{Graph}^{\IndCoh}_{\phi_{\mathcal{Z}},*}(\omega_{\mathcal{Z}}).
\]
这又给出
\[
  \begin{aligned}
    &\Delta^!_{\mathcal{Y}}(\mathrm{id} \times f)^{*,\IndCoh}
      \mathrm{Graph}^{\IndCoh}_{\phi_{\mathcal{Z}} \circ f,*}(\omega_{\mathcal{Y}})\\
    &\qquad\qquad \simeq \Delta^!_{\mathcal{Y}}(\mathrm{id} \times f)^{*,\IndCoh}
      (f \times \mathrm{id})^!\mathrm{Graph}^{\IndCoh}_{\phi_{\mathcal{Z}},*}(\omega_{\mathcal{Z}}).
  \end{aligned}
\]

我们有\footnote{这类交换性是一般事实：见 [12] 命题 7.1.6。不过，在目前的设定中它尤其容易：由 Künneth 公式，我们可以把 $(f \times \mathrm{id})^!$ 写成 $f^! \otimes \mathrm{id}$，把 $(\mathrm{id} \times f)^{*,\IndCoh}$ 写成 $\mathrm{id} \otimes f^{*,\IndCoh}$。}
$(\mathrm{id} \times f)^{*,\IndCoh}(f \times \mathrm{id})^! \simeq (f \times \mathrm{id})^!(\mathrm{id} \times f)^{*,\IndCoh}$，于是可把上式右端改写为
\[
  \begin{aligned}
    &\Delta^!_{\mathcal{Y}}(f \times \mathrm{id})^!(\mathrm{id} \times f)^{*,\IndCoh}
      \mathrm{Graph}^{\IndCoh}_{\phi_{\mathcal{Z}},*}(\omega_{\mathcal{Z}})\\
    &\qquad\qquad = \mathrm{Graph}^{\sigma,!}_f(\mathrm{id} \times f)^{*,\IndCoh}
      \mathrm{Graph}^{\IndCoh}_{\phi_{\mathcal{Z}},*}(\omega_{\mathcal{Z}}).
  \end{aligned}
\]
这里我们用 $\mathrm{Graph}^\sigma_g$ 表示 $(g \times \mathrm{id}) \circ \Delta$，即交换两个笛卡尔因子之后的图像映射。

现在，对纤维积图表
\[
  \begin{tikzcd}[column sep=large, row sep=large]
    \mathcal{Y} \arrow[r, "\mathrm{Graph}^\sigma_f"] \arrow[d, "f"]
      & \mathcal{Z} \times \mathcal{Y} \arrow[d, "\mathrm{id} \times f"] \\
    \mathcal{Z} \arrow[r, "\Delta_{\mathcal{Z}}"]
      & \mathcal{Z} \times \mathcal{Z}
  \end{tikzcd}
  \qquad
  \begin{tikzcd}[column sep=large, row sep=large]
    \mathcal{Z}^\phi \arrow[r, "\varpi_{\mathcal{Z}}"] \arrow[d, "\varpi_{\mathcal{Z}}"]
      & \mathcal{Z} \arrow[d, "\mathrm{Graph}_{\phi_{\mathcal{Z}}}"] \\
    \mathcal{Z} \arrow[r, "\Delta_{\mathcal{Z}}"]
      & \mathcal{Z} \times \mathcal{Z}
  \end{tikzcd}
\]
作基变换，得到同构
\[
  \begin{aligned}
    \mathrm{Graph}^{\sigma,!}_f(\mathrm{id} \times f)^{*,\IndCoh}
      \mathrm{Graph}^{\IndCoh}_{\phi_{\mathcal{Z}},*}(\omega_{\mathcal{Z}})
    &\simeq f^{*,\IndCoh}\Delta^!_{\mathcal{Z}}
      \mathrm{Graph}^{\IndCoh}_{\phi_{\mathcal{Z}},*}(\omega_{\mathcal{Z}})\\
    &\simeq f^{*,\IndCoh}\varpi^{\IndCoh}_{\mathcal{Z},*}(\omega_{\mathcal{Z}^\phi}),
  \end{aligned}
\]
正如所愿。

\begin{remark}
映射 (3.4) 是 $\Gamma(\mathcal{Z}^\phi, \OO_{\mathcal{Z}^\phi})$-线性的。
\end{remark}

\pnumt{范畴框架。}{}\ 我们现在给出处理上述这类构造的一个更概念化的途径。

假设 $F : \mathcal{C} \to \mathcal{D} \in \DGCat_{\mathrm{cont}}$ 是可对偶 DG 范畴之间的映射。假设 $\mathcal{C}$（相应地 $\mathcal{D}$）配有自函子 $\phi_{\mathcal{C}}$（相应地 $\phi_{\mathcal{D}}$），并且：
\begin{itemize}
  \item $F$ 具有连续右伴随 $F^R$。
  \item $F$ 以 lax 方式缠绕 $\phi$，即给定映射（常常是同构）$F \circ \phi_{\mathcal{C}} \to \phi_{\mathcal{D}} \circ F$。
\end{itemize}
于是迹的标准函子性给出与这份数据相伴的典范映射
\[
  \tr_{\mathcal{C}}(\phi_{\mathcal{C}}) \longrightarrow \tr_{\mathcal{D}}(\phi_{\mathcal{D}}) \in \Vect.
\]
确切地说，我们有
\[
  \tr_{\mathcal{C}}(\phi_{\mathcal{C}}) \longrightarrow \tr_{\mathcal{C}}(\phi_{\mathcal{C}} F^R F)
  \simeq \tr_{\mathcal{D}}(F \phi_{\mathcal{C}} F^R) \longrightarrow \tr_{\mathcal{D}}(\phi_{\mathcal{D}} F F^R)
  \longrightarrow \tr_{\mathcal{D}}(\phi_{\mathcal{D}}).
\]
这里我们用了迹的循环性与标准的伴随映射。

\pnumt{推前/拉回再论。}{}\ 首先假设给定如上可表且紧合的 $f : \mathcal{Y} \to \mathcal{Z}$。我们还假设 $\mathcal{Y}$ 与 $\mathcal{Z}$ 是 ind-QCA 的（以确保 $\IndCoh$ 可对偶）。

在 §3.2.5 的设定中，取 $\mathcal{C} = \IndCoh(\mathcal{Y})$、$\mathcal{D} = \IndCoh(\mathcal{Z})$、$F = f^{\IndCoh}_*$、$\phi_{\mathcal{C}} = \phi^{\IndCoh}_{\mathcal{Y},*}$，以及 $\phi_{\mathcal{D}} = \phi^{\IndCoh}_{\mathcal{Z},*}$。

于是范畴形式体系给出典范映射\footnote{这里等号都是标准的同构；例如见 [15] §3.5.3。}：
\[
  \Gamma^{\IndCoh}(\mathcal{Y}^\phi, \omega_{\mathcal{Y}^\phi})
  = \tr_{\IndCoh(\mathcal{Y})}(\phi^{\IndCoh}_{\mathcal{Y},*})
  \longrightarrow \tr_{\IndCoh(\mathcal{Z})}(\phi^{\IndCoh}_{\mathcal{Z},*})
  = \Gamma^{\IndCoh}(\mathcal{Z}^\phi, \omega_{\mathcal{Z}^\phi}).
\]
一次直接的图表追踪表明，这一映射还原出 (3.3)。我们指出，紧合性是 $f^{\IndCoh}_*$ 具有连续右伴随所需要的。

其次，取 $f : \mathcal{Y} \to \mathcal{Z}$ 1-可表且最终余连通。现在取 $\mathcal{C} = \IndCoh(\mathcal{Z})$、$\mathcal{D} = \IndCoh(\mathcal{Y})$、$F = f^{\IndCoh,*}$、$\phi_{\mathcal{C}} = \phi^!_{\mathcal{Z}}$、$\phi_{\mathcal{D}} = \phi^!_{\mathcal{Y}}$。于是范畴形式体系给出典范映射\footnote{注意 $\tr_{\mathcal{C}}(F) = \tr_{\mathcal{C}^\vee}(F^\vee)$，故 $\tr(\phi^!) = \tr(\phi^{\IndCoh}_*)$。}：
\[
  \Gamma^{\IndCoh}(\mathcal{Z}^\phi, \omega_{\mathcal{Z}^\phi})
  = \tr_{\IndCoh(\mathcal{Z})}(\phi^!_{\mathcal{Z}})
  \longrightarrow \tr_{\IndCoh(\mathcal{Y})}(\phi^!_{\mathcal{Y}})
  = \Gamma^{\IndCoh}(\mathcal{Y}^\phi, \omega_{\mathcal{Y}^\phi}).
\]
一次图表追踪表明，这一映射与 (3.4) 一致。

\begin{remark}
第二次图表追踪是例行公事，但不出所料，比第一次要复杂一些。我们在此省略验证。实际上，就我们的目的而言，读者可以把这些范畴构造当作 (3.3) 与 (3.4) 的\emph{定义}，完全忽略 §3.2.3--3.2.4 的材料。我们之所以纳入这些显式构造，只是为了让大家看到构造的具体样子。
\end{remark}
